\documentclass[12pt]{article}
\usepackage[utf8]{inputenc}
\usepackage[T1]{fontenc}
\usepackage[margin=1in]{geometry}
\usepackage{setspace}
\usepackage{xcolor}
\usepackage[colorlinks=true, urlcolor=blue!60!black, linkcolor=blue!60!black]{hyperref}
\setlength{\parindent}{2em}
\setlength{\parskip}{1.25em}
\renewcommand{\baselinestretch}{1.0}
\title{\textit{Brandenburg v. Ohio}, 395 U.S. 444 (1969)}
\author{Case Brief}
\date{}
\begin{document}
\maketitle

\section{Case Metadata}

\begin{itemize}
  \item \textbf{Case Name:} \href{https://supreme.justia.com/cases/federal/us/395/444/}{\emph{Brandenburg v.\ Ohio}}, 395 U.S.\ 444 (1969) (No.\ 492).
  \item \textbf{Court \& Dates:} Supreme Court of the United States; argued February 27, 1969; decided June 9, 1969 (\href{https://www.oyez.org/cases/1968/492}{Oyez}).
  \item \textbf{Vote \& Alignment:} 8--0 unanimous \emph{per curiam} decision reversing the Supreme Court of Ohio (Warren, C.J., and Black, Douglas, Harlan, Brennan, Stewart, White, and Marshall, JJ.; one seat vacant following Justice Fortas's resignation; \href{https://constitutioncenter.org/the-constitution/supreme-court-case-library/brandenburg-v-ohio}{National Constitution Center}).
  \item \textbf{Concurrences \& Dissents:} Black, J., filed a concurring opinion; Douglas, J., filed a concurring opinion. No dissents.
  \item \textbf{Constitutional Provisions \& Statute at Issue:} First and Fourteenth Amendments (Free Speech, Free Press, and Assembly Clauses); Ohio Criminal Syndicalism Statute, Ohio Rev.\ Code Ann.\ \S~2923.13 (enacted 1919).
  \item \textbf{Primary Reference:} \href{https://supreme.justia.com/cases/federal/us/395/444/}{Justia U.S.\ Supreme Court Center}; \href{https://www.law.cornell.edu/wex/brandenburg_test}{Cornell LII Wex}.
\end{itemize}

\section{Facts and Procedural History}

Clarence Brandenburg, a leader of a Ku Klux Klan group, telephoned an announcer-reporter at a Cincinnati television station and invited him to cover a Klan ``rally'' on a farm in Hamilton County, Ohio (\href{https://supreme.justia.com/cases/federal/us/395/444/}{395 U.S.\ at 445}). With the organizers' cooperation, the reporter and a cameraman filmed the gathering, and portions of two films were later broadcast locally and nationally. Only the Klan participants and the news crew were present on the farm.

The first film showed twelve hooded figures---some carrying firearms---gathered around a burning wooden cross, uttering mostly incomprehensible phrases derogatory of Black people and Jews. Another scene on the same film showed Brandenburg, unarmed and in Klan regalia, delivering a speech claiming hundreds of Ohio members and declaring: ``We're not a revengent organization, but if our President, our Congress, our Supreme Court, continues to suppress the white, Caucasian race, it's possible that there might have to be some revengeance taken'' (\href{https://supreme.justia.com/cases/federal/us/395/444/}{395 U.S.\ at 446}). He added that the group planned a July Fourth march on Congress ``four hundred thousand strong,'' splitting toward St.\ Augustine, Florida, and Mississippi. The second film depicted six hooded figures and a similar speech omitting ``revengeance'' while adding derogatory remarks urging that Black people be returned to Africa and Jews to Israel. At trial, Ohio introduced the films, identification testimony, and items shown in the footage (a pistol, rifle, shotgun, ammunition, a Bible, and a red hood).

Ohio convicted Brandenburg under the Ohio Criminal Syndicalism Statute, Ohio Rev.\ Code Ann.\ \S~2923.13, on two counts: (1) advocating ``the duty, necessity, or propriety of crime, sabotage, violence, or unlawful methods of terrorism as a means of accomplishing industrial or political reform,'' and (2) ``voluntarily assembl[ing] with any society, group, or assemblage of persons formed to teach or advocate the doctrines of criminal syndicalism'' (\href{https://supreme.justia.com/cases/federal/us/395/444/}{395 U.S.\ at 444--45}). He was fined \$1,000 and sentenced to one to ten years' imprisonment. An Ohio intermediate appellate court affirmed without opinion, and the Supreme Court of Ohio dismissed his appeal \emph{sua sponte} for lack of a ``substantial constitutional question'' (\href{https://supreme.justia.com/cases/federal/us/395/444/}{395 U.S.\ at 445}). The U.S.\ Supreme Court noted probable jurisdiction and reversed.

\section{Issue(s) Presented}

\begin{enumerate}
  \item Does the Ohio Criminal Syndicalism Statute violate the First and Fourteenth Amendments---facially and as applied---by criminalizing the mere advocacy of force or law violation as a means of political or industrial reform, without requiring that such advocacy be directed to inciting or producing imminent lawless action and likely to do so?
  \item May a State criminalize voluntary assembly with a group formed merely to teach or advocate doctrines of criminal syndicalism without distinguishing mere advocacy from incitement to imminent lawless action?
\end{enumerate}

\section{Holding}

\textbf{Yes} to Issue 1, and \textbf{No} to Issue 2. In a unanimous \emph{per curiam} opinion, the Court reversed Brandenburg's conviction and invalidated Ohio's Criminal Syndicalism Act both on its face and as applied (\href{https://supreme.justia.com/cases/federal/us/395/444/}{395 U.S.\ at 448--49}). The Court held that the First and Fourteenth Amendments prohibit a State from forbidding or proscribing advocacy of the use of force or law violation \emph{except} where such advocacy is: (1) \textbf{directed to inciting or producing imminent lawless action}, and (2) \textbf{likely to incite or produce such action}. The same distinction governs statutes restricting the cognate right of peaceable assembly. The Court expressly overruled \href{https://supreme.justia.com/cases/federal/us/274/357/}{\emph{Whitney v.\ California}}, 274 U.S.\ 357 (1927).

\section{Reasoning}

The \emph{per curiam} opinion situated Ohio's 1919 statute among twenty state and two territorial criminal syndicalism laws enacted between 1917 and 1920 (\href{https://supreme.justia.com/cases/federal/us/395/444/}{395 U.S.\ at 447}). Although the Court had upheld California's substantially similar statute in \href{https://supreme.justia.com/cases/federal/us/274/357/}{\emph{Whitney v.\ California}}, 274 U.S.\ 357 (1927), on the theory that advocating violent political or economic change inherently endangered state security, the Court declared that \emph{Whitney} had been ``thoroughly discredited by later decisions,'' citing \href{https://supreme.justia.com/cases/federal/us/341/494/}{\emph{Dennis v.\ United States}}, 341 U.S.\ 494, 507 (1951), and noting in footnote 2 that \emph{Yates v.\ United States}, 354 U.S.\ 298 (1957), reinforced the requirement of a connection to forcible action.

Distilling those precedents, the Court articulated its landmark two-part incitement test: advocacy of force or illegality is constitutionally protected unless it is both \emph{directed to} inciting or producing \emph{imminent} lawless action (a demanding intent and temporal-immediacy requirement) and \emph{likely} to incite or produce such action (an objective probability requirement). Quoting \href{https://supreme.justia.com/cases/federal/us/367/290/}{\emph{Noto v.\ United States}}, 367 U.S.\ 290, 297--98 (1961), the Court drew a sharp constitutional boundary between ``mere abstract teaching'' of the moral propriety or necessity of violence and ``preparing a group for violent action and steeling it to such action'' (\href{https://supreme.justia.com/cases/federal/us/395/444/}{395 U.S.\ at 447--48}).

Measured against this standard, Ohio's statute failed on its face and as applied. Its text punished anyone who advocated or taught the duty, necessity, or propriety of violence for reform, circulated writings containing such advocacy, or voluntarily assembled with a group formed to teach it. Neither the indictment nor the trial judge's jury instructions---which merely tracked the statutory language---refined the offense to distinguish mere advocacy from incitement to imminent lawless action, and Ohio's courts had never supplied a narrowing construction (\href{https://supreme.justia.com/cases/federal/us/395/444/}{395 U.S.\ at 448--49 \& n.3}). Because the statute swept protected speech and peaceable assembly within its condemnation, the Court struck it down and overruled \emph{Whitney}.

\section{Concurring and Dissenting Opinions}

\begin{itemize}
  \item \textbf{Justice Black (concurring):} In a brief concurrence, Justice Black joined the Court's opinion on the understanding that its citation to \href{https://supreme.justia.com/cases/federal/us/341/494/}{\emph{Dennis v.\ United States}}, 341 U.S.\ 494 (1951), did not endorse the ``clear and present danger'' doctrine on which \emph{Dennis} purported to rely (\href{https://supreme.justia.com/cases/federal/us/395/444/}{395 U.S.\ at 449--50}). He agreed with Justice Douglas that ``clear and present danger'' should have no place in First Amendment interpretation.
  \item \textbf{Justice Douglas (concurring):} Joining the Court's opinion with a ``caveat,'' Justice Douglas traced ``clear and present danger'' from Justice Holmes's World War I opinions in \href{https://supreme.justia.com/cases/federal/us/249/47/}{\emph{Schenck v.\ United States}}, 249 U.S.\ 47 (1919), through its Cold War distortion in \emph{Dennis} (\href{https://supreme.justia.com/cases/federal/us/395/444/}{395 U.S.\ at 450--54}). Douglas argued that the test was easily manipulated by judges wedded to the status quo and was irreconcilable with the First Amendment in peacetime. Rejecting any balancing test, he proposed a categorical line between ``ideas and overt acts'': apart from rare instances where ``speech is brigaded with action'' (such as falsely shouting fire in a crowded theater), pure speech and political advocacy should be immune from prosecution (\href{https://supreme.justia.com/cases/federal/us/395/444/}{395 U.S.\ at 456--57}).
  \item \textbf{Dissents:} None.
\end{itemize}

\section{Key Quotations}

\begin{quote}
``[T]he constitutional guarantees of free speech and free press do not permit a State to forbid or proscribe advocacy of the use of force or of law violation except where such advocacy is directed to inciting or producing imminent lawless action and is likely to incite or produce such action.'' \\
--- \emph{Per curiam}, \href{https://supreme.justia.com/cases/federal/us/395/444/}{395 U.S.\ at 447}.
\end{quote}

\begin{quote}
``[T]he mere abstract teaching . . . of the moral propriety or even moral necessity for a resort to force and violence is not the same as preparing a group for violent action and steeling it to such action.'' \\
--- \emph{Per curiam}, \href{https://supreme.justia.com/cases/federal/us/395/444/}{395 U.S.\ at 447--48} (quoting \href{https://supreme.justia.com/cases/federal/us/367/290/}{\emph{Noto v.\ United States}}, 367 U.S.\ 290, 297--98 (1961)).
\end{quote}

\begin{quote}
``Accordingly, we are here confronted with a statute which, by its own words and as applied, purports to punish mere advocacy and to forbid, on pain of criminal punishment, assembly with others merely to advocate the described type of action. Such a statute falls within the condemnation of the First and Fourteenth Amendments. The contrary teaching of \emph{Whitney v.\ California}, \emph{supra}, cannot be supported, and that decision is therefore overruled.'' \\
--- \emph{Per curiam}, \href{https://supreme.justia.com/cases/federal/us/395/444/}{395 U.S.\ at 449}.
\end{quote}

\begin{quote}
``The line between what is permissible and not subject to control and what may be made impermissible and subject to regulation is the line between ideas and overt acts. The example usually given by those who would punish speech is the case of one who falsely shouts fire in a crowded theatre. This is, however, a classic case where speech is brigaded with action.'' \\
--- Douglas, J., concurring, \href{https://supreme.justia.com/cases/federal/us/395/444/}{395 U.S.\ at 456--57}.
\end{quote}

\section{Doctrinal Significance and Modern Legacy (Through 2026)}

\emph{Brandenburg} stands as the capstone of twentieth-century subversive-advocacy doctrine and one of the most speech-protective rules in comparative constitutional law. By replacing the deferential ``bad tendency'' approach of \emph{Whitney} and the sliding-scale balancing of \emph{Dennis} with a conjunctive requirement of \textbf{specific intent}, \textbf{imminence}, and \textbf{likelihood}, the Court erected an extraordinary shield for dissenting and even hateful political rhetoric.

Subsequent Supreme Court decisions reinforced \emph{Brandenburg}'s rigor. In \href{https://supreme.justia.com/cases/federal/us/414/105/}{\emph{Hess v.\ Indiana}}, 414 U.S.\ 105, 108--09 (1973) (per curiam), the Court reversed the disorderly conduct conviction of an antiwar demonstrator who shouted, ``We'll take the fucking street later (or again),'' holding that advocacy of illegal action at ``some indefinite future time'' fails \emph{Brandenburg}'s imminence prong. In \href{https://supreme.justia.com/cases/federal/us/458/886/}{\emph{NAACP v.\ Claiborne Hardware Co.}}, 458 U.S.\ 886, 927--29 (1982), the Court extended \emph{Brandenburg} to civil tort liability, holding that a civil rights organizer's fiery boycott speeches---warning that violators would have their necks broken---could not support damages absent authorization, ratification, or incitement of imminent violence.

In the 2020s, \emph{Brandenburg} remains the governing benchmark for incitement and a touchstone for mens rea across unprotected speech categories. In \href{https://www.law.cornell.edu/supremecourt/text/22-138}{\emph{Counterman v.\ Colorado}}, 600 U.S.\ \_\_\_ (2023), the Court held that true-threats prosecutions require at least subjective recklessness, expressly contrasting true threats with \emph{Brandenburg} incitement: because incitement lies ``a hair's-breadth away from political `advocacy,''' the First Amendment continues to demand specific intent (purpose or knowledge) for incitement, whereas recklessness suffices for true threats. Justice Sotomayor, concurring in part, cautioned that lowering the true-threats mens rea to recklessness could let prosecutors reframe inflammatory political speech protected under \emph{Brandenburg} and \emph{Claiborne} as reckless true threats.

Through 2026, \emph{Brandenburg}'s boundaries face acute pressure in two arenas: protest-tort litigation and digital-age speech. In \href{https://www.supremecourt.gov/opinions/23pdf/23-373_8njq.pdf}{\emph{Mckesson v.\ Doe}}, 601 U.S.\ \_\_\_ (2024), the Court denied certiorari in a lawsuit seeking to hold a protest organizer liable in negligence for a third party's assault on a police officer, with Justice Sotomayor emphasizing that lower courts must give ``full and fair consideration'' to \emph{Counterman} and \emph{Claiborne}. Yet in \href{https://www.ca5.uscourts.gov/opinions/pub/24/24-30494-CV0.pdf}{\emph{Ford v.\ Mckesson}}, No.\ 24-30494 (5th Cir.\ Mar.\ 19, 2026), \href{https://www.ca5.uscourts.gov/opinions/pub/24/24-30494-CV1.pdf}{reh'g en banc denied (5th Cir.\ July 30, 2026)}, a Fifth Circuit panel allowed the negligence suit to proceed to trial; six judges dissented from denial of en banc review, warning that permitting ``foreseeable violence'' tort liability for civil disobedience ``exhumes'' the bad-tendency test and violates \emph{Brandenburg}. Meanwhile, courts and scholars continue to debate how \emph{Brandenburg}'s face-to-face ``imminence'' and ``intent'' elements apply to viral social media calls to violence and generative AI outputs (\href{https://www.law.cornell.edu/wex/imminent_lawless_action}{Cornell LII Wex}).

\section{Pedagogical Critiques and Exam Hypotheticals}

\subsection*{Key Critiques}
\begin{itemize}
  \item \textbf{Doctrinal Revisionism:} The \emph{per curiam} opinion claimed its test merely synthesized \emph{Dennis} and \emph{Yates}, even though \emph{Dennis} used a gravity-discounted-by-improbability balancing test that dispensed with strict imminence.
  \item \textbf{Vague Scope of ``Imminence'' and ``Lawless Action'':} \emph{Brandenburg} never defined how immediate ``imminent'' must be (seconds, hours, or a scheduled future raid?) or whether ``lawless action'' includes minor misdemeanors like trespass or traffic obstruction, as opposed to serious violence.
  \item \textbf{Digital \& Asymmetric Threats:} Critics argue that \emph{Brandenburg}'s paradigm---a speaker addressing a live crowd---fits poorly with asynchronous online radicalization, viral stochastic incitement, or instructional manuals for crimes.
\end{itemize}

\subsection*{Exam Hypotheticals}
\begin{enumerate}
  \item \textbf{The Scheduled Livestream Raid (Testing Imminence \& Likelihood):} An online political influencer with 500,000 followers livestreams a furious denunciation of a county courthouse and tells viewers, ``Bring crowbars to the courthouse steps at noon tomorrow and smash the doors in so they can't certify the vote.'' No violence occurs during the stream. Under \href{https://supreme.justia.com/cases/federal/us/395/444/}{\emph{Brandenburg}} and \href{https://supreme.justia.com/cases/federal/us/414/105/}{\emph{Hess}}, the prosecution faces a steep imminence hurdle because a 24-hour gap leaves time for counterspeech and police intervention---though courts may scrutinize whether providing a definite time, place, and tools edges toward unprotected solicitation or conspiracy rather than abstract advocacy (\href{https://www.law.cornell.edu/wex/imminent_lawless_action}{Cornell LII Wex}).
  \item \textbf{The ``Opposite Side'' Contrast (Unprotected Incitement vs.\ Protected Advocacy):} Contrast Clarence Brandenburg's filmed speech on a secluded farm (attended only by twelve Klan members and a TV crew, gesturing vaguely to future ``revengeance'') with a rally leader standing outside a besieged municipal building facing an angry, armed crowd of 200 protesters who shouts through a bullhorn, ``The doors are unlocked right now---rush the building and drag the councilmembers out!'' Here, \emph{Brandenburg} flips to the \textbf{opposite side}: the speaker uses explicit words of incitement \emph{directed to} immediate unlawful entry and assault, and given the armed, agitated crowd at the exact target location, lawless action is both \emph{imminent} and objectively \emph{likely}. The First Amendment affords no shield.
\end{enumerate}

\end{document}
